\section{Lecture 27: Diffusion Models}

\subsection{Bridge: Diffusion as a Generative Paradigm}

In the previous lectures, we studied two major generative paradigms:

\begin{itemize}
    \item \textbf{latent-variable likelihood models} such as VAEs,
    \item \textbf{implicit adversarial models} such as GANs.
\end{itemize}

Diffusion models offer a third viewpoint.

\medskip

\noindent
\textbf{Key idea.}  
Instead of generating a complex sample in one step, diffusion models generate by \emph{iterative refinement}:
\begin{itemize}
    \item begin from simple noise,
    \item repeatedly denoise,
    \item gradually recover structure.
\end{itemize}

\medskip

\noindent
This changes the nature of the learning problem. Rather than learning one large jump from latent noise to data, the model learns many small denoising steps.

\medskip

\noindent
\textbf{Why this is attractive.}
\begin{itemize}
    \item generation is broken into simpler local tasks,
    \item training resembles supervised denoising,
    \item optimization is often more stable than adversarial learning,
    \item the model tends to cover diverse modes of the data distribution.
\end{itemize}

\medskip

\noindent
Thus diffusion models can be viewed as generative models built around a simple idea:
\begin{quote}
To learn how to create structure, first learn how to remove noise.
\end{quote}

\subsection{Why Generate by Denoising?}

Direct generation of high-dimensional structured data such as images is difficult because the output space is extremely large and only a tiny region corresponds to realistic data.

\medskip

\noindent
Diffusion models turn this into a sequence of easier problems:
\begin{itemize}
    \item start from a known simple distribution, usually Gaussian noise,
    \item learn how to make a small correction,
    \item repeat many times.
\end{itemize}

\medskip

\noindent
\textbf{Intuition.}
Correcting a slightly noisy image is easier than drawing a coherent image from scratch. If each step is small, the model only has to learn local refinement.

\medskip

\noindent
\textbf{Iterative refinement viewpoint.}
Generation is not a one-shot act but a progressive construction:
\[
x_T \to x_{T-1} \to \cdots \to x_1 \to x_0.
\]

\medskip

\noindent
Each step makes the sample slightly more structured and slightly less noisy.

\subsection{The Forward Noising Process}

Diffusion models begin by defining a \textbf{forward process} that gradually corrupts data.

\medskip

\noindent
Starting from a clean sample $x_0$, we define a Markov chain:
\[
x_0 \to x_1 \to x_2 \to \cdots \to x_T.
\]

At each step, we add a small amount of Gaussian noise:
\[
q(x_t \mid x_{t-1})
=
\mathcal{N}\!\left(
x_t \;\middle|\; \sqrt{1-\beta_t}\,x_{t-1},\; \beta_t I
\right),
\]
where $\beta_t \in (0,1)$ is a prescribed noise schedule.

\medskip

\noindent
Equivalently, we can write
\[
x_t = \sqrt{1-\beta_t}\,x_{t-1} + \sqrt{\beta_t}\,\epsilon_t,
\qquad
\epsilon_t \sim \mathcal{N}(0,I).
\]

\subsection{Interpretation of the Forward Process}

The forward process gradually destroys structure.

\begin{itemize}
    \item for small $t$, $x_t$ is only slightly corrupted,
    \item for intermediate $t$, major structure remains but details are noisy,
    \item for large $t$, $x_t$ becomes nearly indistinguishable from Gaussian noise.
\end{itemize}

\medskip

\noindent
\textbf{Role of the schedule.}
The values $\beta_t$ determine how quickly information is destroyed. Small $\beta_t$ means gentle corruption; larger $\beta_t$ destroys structure faster.

\medskip

\noindent
\textbf{Conceptual point.}
The forward process is \emph{fixed}, not learned. It is chosen by us so that the final state is simple and tractable.

\subsection{Closed-Form Forward Marginals}

A key mathematical fact makes diffusion practical: although the forward process is defined recursively, we can sample $x_t$ directly from $x_0$.

\medskip

\noindent
Define
\[
\alpha_t = 1-\beta_t,
\qquad
\bar{\alpha}_t = \prod_{s=1}^t \alpha_s.
\]

Then:
\[
q(x_t \mid x_0)
=
\mathcal{N}\!\left(
x_t \;\middle|\; \sqrt{\bar{\alpha}_t}\,x_0,\; (1-\bar{\alpha}_t)I
\right).
\]

\medskip

\noindent
Equivalently,
\[
x_t = \sqrt{\bar{\alpha}_t}\,x_0 + \sqrt{1-\bar{\alpha}_t}\,\epsilon,
\qquad
\epsilon \sim \mathcal{N}(0,I).
\]

\medskip

\noindent
\textbf{Why this matters.}
This means we do \emph{not} need to simulate the full forward chain during training. We can:
\begin{itemize}
    \item choose a random time step $t$,
    \item sample one Gaussian noise vector $\epsilon$,
    \item construct $x_t$ directly.
\end{itemize}

\subsection{Intuition Behind the Closed Form}

The formula
\[
x_t = \sqrt{\bar{\alpha}_t}\,x_0 + \sqrt{1-\bar{\alpha}_t}\,\epsilon
\]
shows that the noisy sample is a mixture of:
\begin{itemize}
    \item a shrunk version of the original data,
    \item Gaussian noise with increasing strength.
\end{itemize}

\medskip

\noindent
As $t$ grows:
\begin{itemize}
    \item $\sqrt{\bar{\alpha}_t}$ decreases,
    \item $1-\bar{\alpha}_t$ increases,
    \item the original signal fades while noise dominates.
\end{itemize}

\medskip

\noindent
Thus diffusion constructs a continuum between structured data and pure noise.

\subsection{The Reverse Denoising Process}

Generation is defined by reversing the corruption.

\medskip

\noindent
We posit a learned reverse Markov chain:
\[
p_\theta(x_{0:T})
=
p(x_T)\prod_{t=1}^T p_\theta(x_{t-1}\mid x_t),
\]
where
\[
p(x_T)=\mathcal{N}(0,I).
\]

\medskip

\noindent
So generation begins from
\[
x_T \sim \mathcal{N}(0,I),
\]
then repeatedly samples:
\[
x_{t-1} \sim p_\theta(x_{t-1}\mid x_t),
\]
until we obtain $x_0$.

\medskip

\noindent
\textbf{Interpretation.}
The reverse process is the learned generative model. Each step removes a small amount of noise and restores a small amount of structure.

\subsection{Parameterizing the Reverse Process}

A standard parameterization is Gaussian:
\[
p_\theta(x_{t-1}\mid x_t)
=
\mathcal{N}\!\left(
x_{t-1} \mid \mu_\theta(x_t,t),\; \Sigma_\theta(x_t,t)
\right).
\]

\medskip

\noindent
The neural network receives:
\begin{itemize}
    \item the noisy input $x_t$,
    \item the time step $t$,
\end{itemize}
and predicts quantities describing the reverse transition.

\medskip

\noindent
In practice, modern implementations often predict one of the following:
\begin{itemize}
    \item the noise $\epsilon$,
    \item the clean sample $x_0$,
    \item a score function related to $\nabla_x \log p(x_t)$.
\end{itemize}

These parameterizations are closely related.

\subsection{Why Noise Prediction Is Natural}

Suppose
\[
x_t = \sqrt{\bar{\alpha}_t}\,x_0 + \sqrt{1-\bar{\alpha}_t}\,\epsilon.
\]

If the model can predict the added noise,
\[
\epsilon_\theta(x_t,t) \approx \epsilon,
\]
then it can approximately recover the clean signal:
\[
x_0 \approx
\frac{x_t - \sqrt{1-\bar{\alpha}_t}\,\epsilon_\theta(x_t,t)}
{\sqrt{\bar{\alpha}_t}}.
\]

\medskip

\noindent
\textbf{Why this helps.}
Predicting noise is often easier and more numerically stable than directly predicting the entire previous state.

\medskip

\noindent
\textbf{Conceptual interpretation.}
The network learns:
\begin{quote}
What part of this current sample is noise, and what part is meaningful structure?
\end{quote}

\subsection{ELBO View of Training}

As in VAEs, diffusion training can be motivated by a variational lower bound.

\medskip

\noindent
The forward process defines
\[
q(x_{1:T}\mid x_0)=\prod_{t=1}^T q(x_t\mid x_{t-1}),
\]
and the reverse model defines
\[
p_\theta(x_{0:T})=p(x_T)\prod_{t=1}^T p_\theta(x_{t-1}\mid x_t).
\]

\medskip

\noindent
We want to maximize $\log p_\theta(x_0)$. Using variational reasoning gives the ELBO:
\[
\mathbb{E}_q
\left[
\log p_\theta(x_{0:T}) - \log q(x_{1:T}\mid x_0)
\right].
\]

\medskip

\noindent
After algebraic simplification, this leads to a practical objective that can be written in a much simpler form.

\subsection{Usable Training Objective}

The most common training loss is:
\[
\mathcal{L}_{\text{diff}}
=
\mathbb{E}_{t,x_0,\epsilon}
\left[
\left\|
\epsilon - \epsilon_\theta(x_t,t)
\right\|^2
\right],
\]
where
\[
x_t = \sqrt{\bar{\alpha}_t}x_0 + \sqrt{1-\bar{\alpha}_t}\epsilon,
\qquad
\epsilon \sim \mathcal{N}(0,I).
\]

\medskip

\noindent
\textbf{Training procedure.}
For each minibatch:
\begin{enumerate}
    \item sample clean data $x_0$,
    \item sample a time index $t$,
    \item sample Gaussian noise $\epsilon$,
    \item construct $x_t$ directly,
    \item train the network to predict $\epsilon$ from $(x_t,t)$.
\end{enumerate}

\medskip

\noindent
\textbf{Important point.}
Although the diffusion model defines a complex probabilistic chain, the practical loss is just a denoising regression problem.

\subsection{Why the Objective Is Stable}

The diffusion objective has a major practical advantage: it resembles supervised learning.

\medskip

\noindent
\textbf{The target is known.}
During training, we ourselves sample the noise $\epsilon$, so the model has an explicit regression target.

\medskip

\noindent
\textbf{Contrast with GANs.}
In GANs, the learning signal is produced by an adversarial discriminator and can become unstable. In diffusion, the model learns from direct denoising targets.

\medskip

\noindent
\textbf{Result.}
This often leads to:
\begin{itemize}
    \item smoother optimization,
    \item less mode collapse,
    \item more predictable training behavior.
\end{itemize}

\subsection{Score Function Intuition}

A central concept in diffusion theory is the \textbf{score function}:
\[
s(x)=\nabla_x \log p(x).
\]

\medskip

\noindent
This vector points toward directions in which the log-density increases most rapidly.

\medskip

\noindent
\textbf{Interpretation.}
If a point lies in a low-probability region, the score points toward more likely regions. Thus score information can guide a sample back toward the data manifold.

\medskip

\noindent
In noisy settings, the relevant object is the score of the noisy distribution at time $t$.

\subsection{Noise Prediction and Score Prediction}

The noise-prediction view and the score view are closely connected.

\medskip

\noindent
When
\[
x_t = \sqrt{\bar{\alpha}_t}x_0 + \sqrt{1-\bar{\alpha}_t}\epsilon,
\]
predicting the noise is equivalent, up to scaling, to predicting a score:
\[
\epsilon_\theta(x_t,t)
\;\Longleftrightarrow\;
-\nabla_x \log q(x_t).
\]

\medskip

\noindent
\textbf{Intuition.}
If the model knows which direction corresponds to ``noise,'' then it also knows which direction moves toward higher data density.

\medskip

\noindent
Thus denoising can be interpreted as moving samples toward regions of higher probability.

\subsection{Reverse Denoising as Probability-Guided Refinement}

Each reverse step can be viewed as performing two related actions:

\begin{itemize}
    \item \textbf{remove estimated noise,}
    \item \textbf{move toward more likely structure.}
\end{itemize}

\medskip

\noindent
This gives diffusion a geometric interpretation:
\begin{quote}
Generation proceeds by repeatedly nudging noise toward the data manifold.
\end{quote}

\medskip

\noindent
This viewpoint connects diffusion models to:
\begin{itemize}
    \item score matching,
    \item stochastic dynamics,
    \item iterative optimization.
\end{itemize}

\subsection{Iterative Refinement as a General Paradigm}

A useful way to understand diffusion is not merely as a probability model, but as a broad computational strategy.

\medskip

\noindent
\textbf{Iterative refinement principle.}
Instead of solving a hard generation task in one step, solve it as a sequence of easier corrections:
\[
x_{t-1}=x_t+\text{small update}.
\]

\medskip

\noindent
This resembles:
\begin{itemize}
    \item residual learning,
    \item gradient-based optimization,
    \item classical iterative reconstruction algorithms.
\end{itemize}

\medskip

\noindent
\textbf{Thought experiment.}
Start from random pixels:
\begin{itemize}
    \item first only coarse structure appears,
    \item then object shapes emerge,
    \item then texture and detail sharpen,
    \item finally the sample becomes realistic.
\end{itemize}

\medskip

\noindent
This layered emergence of structure is one reason diffusion models are so compelling conceptually.

\subsection{Langevin Dynamics Connection}

Diffusion also has a useful connection to Langevin dynamics.

\medskip

\noindent
A simplified Langevin-style update takes the form
\[
x_{t-1}
=
x_t + \eta \nabla_x \log p(x_t) + \sqrt{2\eta}\,\xi,
\]
where $\xi$ is Gaussian noise.

\medskip

\noindent
\textbf{Interpretation.}
\begin{itemize}
    \item the gradient term pushes toward high-density regions,
    \item the noise term maintains stochastic exploration.
\end{itemize}

\medskip

\noindent
Diffusion sampling can be viewed as a learned, time-dependent version of this idea: stochastic updates guided by an estimated score or denoising direction.

\subsection{Continuous-Time View}

If the number of diffusion steps becomes very large, the process approaches a continuous-time stochastic differential equation (SDE):
\[
dx = f(x,t)\,dt + g(t)\,dW_t.
\]

\medskip

\noindent
The reverse generative process is then another SDE that removes noise over time.

\medskip

\noindent
\textbf{Why this matters conceptually.}
It places diffusion at the intersection of:
\begin{itemize}
    \item probability,
    \item dynamical systems,
    \item stochastic processes.
\end{itemize}

\medskip

\noindent
For this lecture, however, the discrete-time viewpoint is enough to understand the main mechanism and training objective.

\subsection{DDPM: Stochastic Sampling}

The classic diffusion sampler is \textbf{DDPM} (Denoising Diffusion Probabilistic Model).

\medskip

\noindent
Its main features are:
\begin{itemize}
    \item reverse sampling uses a learned probabilistic transition,
    \item stochastic noise is injected during sampling,
    \item many small steps are taken from $x_T$ to $x_0$.
\end{itemize}

\medskip

\noindent
\textbf{Interpretation.}
Because randomness remains in the reverse process, repeated runs can generate different outputs even from similar initial conditions.

\medskip

\noindent
\textbf{Benefit.}
This stochasticity supports diversity.

\medskip

\noindent
\textbf{Cost.}
Sampling can be slow because many denoising steps are required.

\subsection{DDIM: Deterministic Sampling}

A widely used alternative is \textbf{DDIM} (Denoising Diffusion Implicit Model).

\medskip

\noindent
DDIM modifies the reverse sampling procedure so that it can follow a deterministic trajectory.

\medskip

\noindent
\textbf{Main features.}
\begin{itemize}
    \item reduced or removed stochasticity,
    \item a more direct trajectory through latent-noise space,
    \item often fewer sampling steps.
\end{itemize}

\medskip

\noindent
\textbf{Interpretation.}
Instead of wandering stochastically while denoising, the model follows a more controlled path.

\medskip

\noindent
\textbf{Benefit.}
Sampling is often faster and more controllable.

\medskip

\noindent
\textbf{Tradeoff.}
Less stochasticity can reduce diversity, though in practice DDIM often performs very well.

\subsection{DDPM vs DDIM Intuition}

The difference can be summarized intuitively as follows.

\medskip

\noindent
\textbf{DDPM.}
\begin{itemize}
    \item more probabilistic,
    \item more stochastic,
    \item often more faithful to the original generative process.
\end{itemize}

\medskip

\noindent
\textbf{DDIM.}
\begin{itemize}
    \item more deterministic,
    \item often faster,
    \item often easier to steer or accelerate.
\end{itemize}

\medskip

\noindent
\textbf{Conceptual lesson.}
Once a denoising model is learned, there is freedom in how we traverse the reverse path from noise to data. Different samplers trade off diversity, speed, and controllability.

\subsection{Why Diffusion Models Work Well in Practice}

Diffusion models have become highly successful for several reasons.

\begin{itemize}
    \item \textbf{Stable training:} the objective is a denoising regression loss.
    \item \textbf{Good mode coverage:} they are typically less prone to collapse than GANs.
    \item \textbf{Iterative refinement:} complex generation is decomposed into simpler steps.
    \item \textbf{Flexibility:} they adapt well to images, audio, video, and multimodal generation.
\end{itemize}

\medskip

\noindent
\textbf{Main weakness.}
Sampling is often slower than one-shot generation, because many reverse steps may be needed.

\medskip

\noindent
Thus diffusion often trades speed for robustness and sample quality.

\subsection{Comparison with Earlier Generative Models}

\textbf{Compared with VAEs.}
\begin{itemize}
    \item diffusion usually produces sharper samples,
    \item VAEs provide a more direct latent-variable inference framework,
    \item both are grounded in probabilistic thinking.
\end{itemize}

\medskip

\noindent
\textbf{Compared with GANs.}
\begin{itemize}
    \item diffusion is usually more stable to train,
    \item GANs can sample faster,
    \item diffusion is less vulnerable to mode collapse,
    \item GANs rely on adversarial games, diffusion relies on denoising supervision.
\end{itemize}

\medskip

\noindent
\textbf{Big-picture view.}
Diffusion models show that high-quality generation does not require either exact likelihood tractability or adversarial competition. Iterative denoising is itself a powerful generative principle.

\subsection{Summary}

\textbf{Main points of this lecture.}

\begin{itemize}
    \item Diffusion models generate by reversing a gradual corruption process.
    \item The forward process is a Gaussian Markov chain:
    \[
    q(x_t \mid x_{t-1})
    =
    \mathcal{N}\!\left(
    x_t \;\middle|\; \sqrt{1-\beta_t}\,x_{t-1},\; \beta_t I
    \right).
    \]
    \item A key closed-form identity is
    \[
    x_t
    =
    \sqrt{\bar{\alpha}_t}\,x_0
    +
    \sqrt{1-\bar{\alpha}_t}\,\epsilon,
    \qquad
    \epsilon \sim \mathcal{N}(0,I).
    \]
    \item The reverse model learns
    \[
    p_\theta(x_{t-1}\mid x_t),
    \]
    which performs denoising step by step.
    \item In practice, training often reduces to noise prediction:
    \[
    \mathbb{E}_{t,x_0,\epsilon}
    \left[
    \|\epsilon-\epsilon_\theta(x_t,t)\|^2
    \right].
    \]
    \item Noise prediction is closely related to score estimation.
    \item Diffusion can be understood as iterative refinement toward higher-density regions.
    \item DDPM uses stochastic reverse sampling, while DDIM uses more deterministic and often faster trajectories.
\end{itemize}

\medskip

\noindent
\textbf{Final insight.}  
Diffusion models succeed because they transform generation from a single difficult act into a long sequence of simple denoising corrections. This makes them both mathematically elegant and practically powerful.